\subsection{平面上的正相似变换。}

回忆（§ 6，第 5 小节），E 的正相似变换是向量空间 E 的一个自同构 u，使得 \(\Phi(u(x), u(y)) = (\det u)\Phi(x, y)\) （对 E 中一切 \(x, y\) ）。

命题 1. --- 设 A(Φ) 是 \(\mathfrak{L}_{\mathbf{A}}(\mathbf{E})\) 的由 E 的正相似变换生成的子代数。

\begin{enumerate}
\def\labelenumi{\alph{enumi})}
\item
  正相似变换就是 A(Φ) 的可逆元素。代数 A(Φ) 是 A 上次数为 2 的交换代数。当 E 不含非零迷向向量时，A(Φ) 是一个域，是 A 的一个二次扩张；否则它是同构于 A 的两个域的直积。
\item
  设 \((e_{1}, e_{2})\) 是 E 的一个正交基；令 \(\alpha_{i} = \Phi(e_{i}, e_{i})\) （i = 1, 2），并令 \(\delta = -\alpha_{2}/\alpha_{1}\) 。A( \(\Phi\) ) 的元素关于这个基的矩阵是形如 \(\begin{pmatrix} a & \delta b \\ b & a \end{pmatrix}\) 的矩阵，其中 \(a \in A\) ，\(b \in A\) 。
\item
  空间 E 是一个自由循环 A(Φ)-模，由任一非迷向向量生成。
\end{enumerate}

事实上，引入 E 上一个辅助的交错双线性型 B ≠ 0；于是 B 非退化。因此存在 E 的一个自同态 w，使得 \(\Phi(x, y) = \mathrm{B}(w(x), y)\) （对 E 中一切 x、y）。对 E 的每个可逆自同态 u，我们有

\[
\begin{array}{r l} \Phi (u (x), u (y)) & = \mathrm{B} (w u (x), u (y)) = \mathrm{B} (u u ^ {- 1} w u (x), u (y)) \\ & = (\det u) \mathrm{B} (u ^ {- 1} w u (x), y); \end{array}
\]

要使 u 是正相似变换，必须且只需

\[
(\det u) \mathrm{B} (u ^ {- 1} w u (x), y) = (\det u) \Phi (x, y) = (\det u) \mathrm{B} (w (x), y)
\]

（对 E 中一切 \(x, y\)）；由于 det \(u \neq 0\) 且 B 非退化，这等价于 \(u^{-1}wu = w\) ，也等价于 \(uw = wu\) 。取 B 为其矩阵 \(S\) 关于 \((e_1, e_2)\) 为 \(\begin{pmatrix} 0 & 1 \\ -1 & 0 \end{pmatrix}\) 的那个交错双线性型；以 \(R\) 与 \(W\) 分别记 \(\Phi\) 与 \(w\) 关于这个基的矩阵，则关系 \(\Phi(x, y) = \mathrm{B}(w(x), y)\) 依 § 1 第 10 小节的公式 (47) 写成 \(R = ^t W.S\) ；把它明写出来即得 \(W = \begin{pmatrix} 0 & \alpha_2 \\ -\alpha_1 & 0 \end{pmatrix}\) 。若以 \(\begin{pmatrix} a & c \\ b & d \end{pmatrix}\) 记 \(u\) 关于 \((e_1, e_2)\) 的矩阵，则关系 \(uw = wu\) 就等价于关系 \(b\alpha_2 = -c\alpha_1\) ，\(a\alpha_2 = d\alpha_2\) ，\(a\alpha_1 = d\alpha_1\) ，也就是 \(a = d\) 与 \(c = \delta b\) ；这证明了正相似变换的矩阵是形如 \(\begin{pmatrix} a & \delta b \\ b & a \end{pmatrix}(a, b\) 在 A 中)的可逆矩阵。

现在，其矩阵形如 \(\begin{pmatrix} a & \delta b \\ b & a \end{pmatrix}\) （a, b 在 A 中）的 E 的自同态构成 \(\mathfrak{L}_{\mathrm{A}}(\mathrm{E})\) 的一个维数为 2 的向量子空间，它由 1 与自同态 w 生成；由于 \(w^{2}\) 是比为 \(-\alpha_{1}\alpha_{2}\) 的位似，这个子空间就是子代数 \(\mathrm{A}(\Phi)\) ，即 \(\mathfrak{L}_{\mathrm{A}}(\mathrm{E})\) 中由正相似变换生成的那个子代数。正相似变换是 \(\mathrm{A}(\Phi)\) 的可逆元素，也就是其矩阵满足 \(a^{2}-\delta b^{2}\neq0\) 的那些元素。代数 \(\mathrm{A}(\Phi)\) 是交换的这一事实是显然的。对它应用第二章 § 7 第 7 小节的结果：若 δ 不是 A 中的平方，即 E 的非零向量都不是迷向的，则 \(\mathrm{A}(\Phi)\) 是一个域；反之，若 δ 是 A 中的平方，即 E 含有非零迷向向量，则 \(\mathrm{A}(\Phi)\) 是同构于 A 的两个域的直和。这就证明了 a) 与 b)。

最后，E 的任一非迷向向量都可以取作

某个正交基的第一个向量 \(e_1\) ，设该正交基为 \((e_1, e_2)\) （ \(\S 6, n^0 1\) ）；因此它的像 \(u(e_1)\) （ \(u\) 为 A( \(\Phi\) ) 的元素）都是形如 \(ae_1 + be_2\) （ \(a, b\) 在 A 中）的向量，也就是 E 中所有的向量，因为所有矩阵 \(\begin{pmatrix} a & \delta b \\ b & a \end{pmatrix}\) 都是 A( \(\Phi\) ) 的元素的矩阵。换言之，E 是一个循环 A( \(\Phi\) )-模，由任一非迷向向量生成。此外我们看到，它是单生成的自由 A( \(\Phi\) -模，因为 \(u(e_1) = ae_1 + be_2 = 0\) 蕴含 \(a = b = 0\) ，从而 \(u = 0\) 。这就证明了 \(c\) )。

注. --- 1) 设 \(\vartheta\) 是矩阵为 \(\begin{pmatrix}0 & \delta \\ 1 & 0\end{pmatrix}\) （关于 \((e_1, e_2)\) ）的相似变换，下面把它记作 v；命题 1 的证明中引入的相似变换 \(w\) 等于 \(-\alpha_1 v\) ；我们有 \(v^2 = \delta\) 。正相似变换 \(u = a + b v\) （ \(a, b\) 在 A 中）的乘子等于它的矩阵 \(\begin{pmatrix}a & \delta b \\ b & a\end{pmatrix}\) 的行列式，即 \(a^2 - \delta b^2 = (a + b v)(a - b v) = u.\bar{u}\) ，其中以 \(\bar{u}\) 记 \(u\) 在代数 A(Φ) 中的共轭元（第二章 § 7，第 7 小节）；换言之，\(u\) 的乘子就是 \(u\) 在代数 A(Φ) 中的范数 N( \(u\) ) （出处同前）。特别地，要使正相似变换 \(u\) 是一个旋转，必须且只需 N( \(u\) ) = 1；要使 \(u\) 是位似，必须且只需 \(u \in A^*\) 。

\begin{enumerate}
\def\labelenumi{\arabic{enumi})}
\setcounter{enumi}{1}
\item
  其平方是位似且形如 \(u = a + bv\) （ \(a, b\) 在 A 中，\(b \neq 0\) ）的正相似变换只有位似以及标量倍 \(bv\) （ \(v\) 的），因为 \((a + bv)^{2} = (a^{2} + \delta b^{2}) + 2abv\) 。后者正是把每个向量变成正交向量的向量空间 E 的那些自同构；事实上，这样的自同构的矩阵必然形如 \(\begin{pmatrix} 0 & c \\ d & 0 \end{pmatrix} (c, d\) 在 A 中)，而把向量 \(\lambda e_{1} + \mu e_{2}\) 变成正交向量的条件这时就写成 \(\lambda \mu(d\alpha_{2} + c\alpha_{1}) = 0\) 。
\item
  容易验证，对 E 中的 \(x\) 、\(y\) 及 \(u \in A(\Phi)\) ，有 \(\Phi(u(x), y) =: \Phi(x, \overline{u}(y))\) 。于是，正相似变换 \(u\) 的伴随自同态是正相似变换 \(\overline{u}\) ，即 \(u\) 在 A( \(\Phi\) ) 中的共轭元。
\item
  由于 E 的每个逆相似变换都是一个正相似变换与关于 \(Ae_{1}\) 的对称（变换）的乘积，逆相似变换关于 \((e_{1}, e_{2})\) 的矩阵是形如 \(\begin{pmatrix} a & -\delta b \\ b & -a \end{pmatrix}\) 的矩阵。
\end{enumerate}

此后我们以 S 表示 E 的相似变换群，以 \(S^{+}\) 表示正相似变换群，以 H 表示非零位似的群，并以 \(O^{+}\) 表示旋转群。回忆我们有 \(H \subset S^{+}\) （§ 6，第 5 小节）。

推论 1. --- 正相似变换的群 S \(^{+}\) 是交换的。对 E 中任意非迷向向量 x、y，存在唯一的正相似变换 u 使 y = u(x)。

第一个断言由代数 A(Φ) 是交换的这一事实得出。由于 E 是由 x（相应地，y）生成的自由单生成 A(Φ)-模，存在 A(Φ) 的唯一元素 u（相应地，u′）使得 \(y = u(x)\) （相应地，\(x = u'(y)\) ）；于是 \(x = u(u'(x))\) ，且 \(uu'\) 是恒等映射；这表明 u 可逆，因此是一个正相似变换。

推论 2. --- 旋转群 O \(^{+}\) 是交换的。对 E 中满足 Φ(x, x) = Φ(y, y) ≠ 0 的任意向量 x、y，存在唯一的旋转 u 使 y = u(x)。

第一个断言由推论 1 得出。它还表明，存在正相似变换 \(u\) 且恰好一个使 \(y = u(x)\) ；由于 \(\Phi(u(x), u(x)) = \Phi(x, x)\) ，\(u\) 的乘子等于 1，因此 \(u\) 是一个旋转。

推论 3. --- 群 S \(^{+}\) /H 是交换的。它作用于 E 的非迷向直线的集合上。无论 E 的非迷向直线 D、D′ 如何，都存在且只存在 S \(^{+}\) /H 的一个元素 φ 使 D' = φ(D)。

这由推论 1 以及 H 使 E 的每条直线整体不变这一事实得出。

命题 2. --- 从 O \(^{+}\) 到 S \(^{+}\) /H 的典范同态的核是 \{1, -1\}。

事实上，1 与 -1 是 \(\mathbf{H} \cap \mathbf{O}^{+}\) 的仅有的元素。

命题 3. --- 同态 \(u \to u/\overline{u} = u^{2}/\mathrm{N}(u)\) （\(\mathrm{S}^{+}\) 到自身的）以 H 为核，并通过取商定义一个从 \(\mathrm{S}^{+}/\mathrm{H}\) 到 \(\mathrm{O}^{+}\) 上的同构。

事实上，关系 \(u/\overline{u} = 1\) 等价于 \(u = \overline{u}\) ，也就是 \(u \in A^* = H\) 。于是 H 是 \(u \to u/\overline{u}\) 的核。由于 \(N(u/\overline{u}) = 1\) ，\(u/\overline{u}\) 是一个旋转（注 1）。剩下要证明每个元素 \(v\) （\(O^+\) 中的）都形如 \(u/\overline{u}\) （ \(u \in S^+\) ）。若 \(1 + v\) 可逆，则可取 \(u = 1 + v\) ，因为关系 \(N(v) = v\bar{v} = 1\) 蕴含 \(1 + v = v(1 + \bar{v})\) 。否则有 \(N(1 + v) = (1 + v)(1 + \bar{v}) = 0\) ，也就是说，令 \(v = a + bw\) （ \(a \in A\) ，\(b \in A\) ，\(w^2 = \delta \in A\) ）后，\(2(1 + a) = 0\) ，从而 \(a = -1\) ；但关系 \(a = -1\) 与 \(N(v) = a^2 - \delta b^2 = 1\) 蕴含 \(b = 0\) ，从而 \(v = -1\) ；由于 \(\bar{w} = -w\) ，在这种情形取 \(u = w\) 即可。

当 A(Φ) 是域时，命题 3 是希尔伯特定理的一个特殊情形（第五章 § 11，第 5 小节，定理 3）。

推论. — 设 \(i: \mathrm{O}^{+} \to \mathrm{S}^{+}/\mathrm{H}\) 与 \(d: \mathrm{S}^{+}/\mathrm{H} \to \mathrm{O}^{+}\) 是命题 2 与 3 中定义的同态，并把交换群 \(\mathrm{O}^{+}\) 与 \(\mathrm{S}^{+}/\mathrm{H}\) 写成加法；则有 \(d(i(\theta)) = 2\theta\) （对 \(\theta \in \mathrm{O}^{+}\) ），且有 \(i(d(\varphi)) = 2\varphi\) （对 \(\varphi \in \mathrm{S}^{+}/\mathrm{H}\) ）。

事实上，若 \(\bar{\nu}\) 是一个旋转，则 \(\bar{\nu} = \nu^{-1}\) ，于是 \(d(i(\nu)) = \nu/\bar{\nu} = \nu^{2}\) 。另一方面，若 \(\varphi \in S^{+}/H\) ，则 \(\varphi\) 是某个正相似变换 \(u\) 的模 H 的类，而 \(d(\varphi) = u/\bar{u} = u^{2}/N(u)\) 模 H 同余于 \(u^{2}\) ；由此得第二个公式。

